\documentclass[10pt,a4paper]{amsart}
\usepackage[margin=19mm]{geometry}
\usepackage{amsmath,amssymb,mathtools}
\usepackage[scr=rsfs,cal=euler]{mathalfa}
\usepackage{xfrac}
\usepackage[unicode,hidelinks,bookmarks=false]{hyperref}
\newcommand{\ie}{i.e.\ }
\newcommand{\cf}{cf.\ }
\newcommand{\resp}{respectively\ }
\newcommand{\Subsection}{\subsection}
\providecommand{\nobreakdash}{\nobreak\hbox{-}}
\setlength{\emergencystretch}{2em}
\multlinegap0pt
\usepackage{mathtools}
\usepackage{amssymb}
\usepackage[shortlabels]{enumitem}
\newtheorem{corollary}{Corollary}
\title{Capparelli's partition theorem as part of an infinite hierarchy:\\ Combinatorial and Weighted Words extensions of recent work}
\author{Yazan Alamoudi, Krishnaswami Alladi}
\date{Source: arXiv:2606.12359v1 (CC BY 4.0)}
\begin{document}
\maketitle

\noindent\emph{Source and licence.} Capparelli's partition theorem as part of an infinite hierarchy:\\ Combinatorial and Weighted Words extensions of recent work. Version arXiv:2606.12359v1 (CC BY 4.0), distributed by arXiv under the Creative Commons Attribution 4.0 International licence, http://creativecommons.org/licenses/by/4.0/, as recorded in the arXiv licence metadata for this exact version and shown on the abstract page https://arxiv.org/abs/2606.12359v1. Full licence notice: \texttt{licenses/arxiv-2606.12359-CC-BY-4.0.txt}.\par\smallskip

\noindent\textit{Editorial scope.} Counted unit: the article's corollary corollary (source0.tex lines 832--836). The article's complete original proof of it is delivered with it (source0.tex lines 837--848, one \texttt{proof} environment of 12 lines). No mathematics is written, repaired or re-derived: the statement and the proof are the article's own lines, in the article's own notation, and no retained line is substituted or re-flowed.  No further source-local context is needed: every cross-reference of every retained line resolves inside the two delivered spans.  Nothing was omitted from any retained span, so the package carries no omission record.  No retained line contains a citation, so the wrapper carries no bibliography and no bibliography entry.  No retained line contains a figure environment, an image-inclusion command or drawing code, so no image is delivered and the package declares no asset dependency.  Editorial formatting only, in addition: an AMS \texttt{amsart} preamble that restates the article's own declarations and macros used by the retained lines, namely the environments \texttt{corollary} on separate counters, together with the standard packages those lines need.  The retained environments print in this document as Corollary 1 in order of appearance, whereas the article prints the same items with the numbering of its own full document.  The complete native source of the article as distributed in the arXiv e-print for this exact version is bundled unchanged as \texttt{sources/202619/source0.tex}.
\begin{corollary}
Let $H_c(i,j,k)_x$ be the generating function for minimal partitions whose smallest part is a $c$-part and smallest non-$c$ part an $x$-part. Then, we have the following. 
\[H_c(i,j,k)_a=q^{2T_i+2T_j+T_k+(i+j)k+2j}{i+j+k-1\brack i+j,k-1}_q{i+j-1\brack i-1,j}_{q^2}\]
\[H_c(i,j,k)_b=q^{2T_i+2T_j+T_k+(i+j)(k+1)}{i+j+k-1\brack i+j,k-1}_q{i+j-1\brack i,j-1}_{q^2}\]
\end{corollary}

\begin{proof}
The proof is the same as the last proof we did for Theorem 3. The only difference is that when we are inducting, the previous cases are generating functions for minimal partitions and not just mock-minimal partitions. For the $a$-parts, we are increasing every part by one every time we add a $c$ part. Like before, this amounts to 
\begin{align*}
    H_c(i,j,k)_{a}=&\sum_{i'=1}^k q^{i'(i+j+k)-T_{i'-1}}H_{a}(i,j,k-i')\\&=q^{2T_i+2T_j+T_k+(i+j)k+2j}{i+j-1\brack i-1,j}_{q^2}\times \sum_{i'=1}^kq^{k-i'}{i+j-1+k-i'\brack i+j-1,k-i'}_{q}\\
    &=q^{2T_i+2T_j+T_k+(i+j)k+2j}{i+j+k-1\brack i+j,k-1}_q{i+j-1\brack i-1,j}_{q^2}.
\end{align*}
For the $b$-parts, it is the same except that the \textbf{first time you add a $c$-part, you increase every part by 2} (this is the combinatorial reason for the extra $q^{i+j}$ factor). After, we are increasing every part by one every time we add a $c$ part. This amounts to 
\begin{align*}
    H_c(i,j,k)_{b}=& \sum_{i'=1}^k q^{i+j+k+i'(i+j+k)-T_{i'-1}}H_{b}(i,j,k-i')\\&=q^{2T_i+2T_j+T_k+(i+j)(k+1)}{i+j-1\brack i-1,j}_{q^2}\times \sum_{i'=1}^kq^{k-i'}{i+j-1+k-i'\brack i+j-1,k-i'}_{q}\\
    &=q^{2T_i+2T_j+T_k+(i+j)(k+1)}{i+j+k-1\brack i+j,k-1}_q{i+j-1\brack i,j-1}_{q^2}
\end{align*}
\end{proof}

\end{document}
